%================================================================
\chapter*{Introduction}
\addcontentsline{toc}{chapter}{Introduction}
\markboth{Introduction}{Introduction}
%================================================================

This document is a companion reference for participants in the
\ProjectName{} program. It covers the software environment, programming
foundations, and data challenge materials needed before and during the
summer bootcamp. It is \emph{not} intended to be read cover to cover.
Think of it as a reference manual: use the routing table on the next
page to identify which chapters apply to your background, complete
those chapters before the first session, and return to the others as
needed during the program.

\noindent
This document is version-locked to the \TheDate{} cohort and is
planned for obsolescence on \TheObsolescenceDate{}. After that date,
retrieve the current edition from the program repository at
\url{\ProjectGit}. The companion curriculum document,
\textit{\ProjectName{} 2026 Curriculum}, covers the seven program
units delivered during the bootcamp; this document covers the
prerequisites and data challenge only.

\begin{infobox}%[title=This document is a reference, not a reading list]
You are \textbf{not} expected to read this document from cover to cover.
Identify the chapters relevant to your background using the routing
table on the next page, complete those before the first session, and
return to the remaining chapters only as the need arises during the
program. Most participants will skip at least two or three chapters
entirely.
\end{infobox}

%================================================================
\section*{Conventions Used in This Document}
%================================================================

\textbf{Inline Q\&A prompts.} Throughout the chapters, red prompts
marked \textcolor{HW}{\textbf{Record your explanation as answer
\#k.}} appear at key points. These are classroom discussion prompts
that do not require a written response here; use the fillable form in
Chapter~\ref{ch:assessment} to record responses to the curated subset
of questions selected for the progress record.

\noindent
\textbf{Callout boxes.} Four types of callout box appear in the text:

\begin{itemize}
    \item \textbf{Information} (teal border): reference material and
          neutral context.
    \item \textbf{Tip} (blue border): practical guidance and
          shortcuts.
    \item \textbf{Warning} (magenta border): critical cautions
          that, if ignored, will cause errors or data loss.
    \item \textbf{Note} (orange border): items requiring attention
          but not critical.
\end{itemize}

\noindent
\textbf{Code listings.} All commands intended to be typed in a
terminal are shown in monospace code blocks. Windows commands use a
dark background; Linux/macOS commands use a lighter background.
Commands that work identically on all platforms use the neutral
listing style.

\noindent
\textbf{Cross-references.} Blue underlined text is a clickable
hyperlink. References to chapters, sections, figures, and tables are
also clickable in the PDF.

%================================================================
\section*{Chapter Dependencies}
%================================================================

      \renewcommand{\thefigure}{\arabic{figure}}
      \setcounter{figure}{0}
      Figure~\ref{fig:dependencies} shows the dependency structure of the
      chapters. A solid arrow indicates a hard dependency: the target chapter
      requires the source to be completed first. A dashed arrow indicates
      recommended order only.

      \begin{figure}[htbp]
      \centering
      \begin{tikzpicture}[
      node distance = 0.55cm and 1.8cm,
      every node/.style = {
            rectangle, rounded corners=4pt,
            minimum width=3.8cm, minimum height=0.7cm,
            align=center, font=\small\sffamily,
            text=white
      },
      harrow/.style  = {-Latex, thick},
      darrow/.style  = {-Latex, dashed, thick}
      ]



% ---- Nodes ----
% Column A: main prerequisite spine
\node[fill=colorEssential]  (env)  {Ch.1 Environment Setup};
\node[fill=colorDesirable,  below=of env]  (cli)  {Ch.2 CLI \& Scripting};
\node[fill=colorDesirable,  below=of cli]  (prog) {Ch.3 Intro to Programming};
\node[fill=colorEssential,  below=of prog] (pyenv){Ch.4 Python Environments};
\node[fill=colorDesirable,  below=of pyenv](numop){Ch.5 Numerical Operations};
\node[fill=colorDesirable,  below=of numop](git)  {Ch.7 Git \& GitHub};
\node[fill=colorOptional,   below=of git]  (llm)  {Ch.8 LLM Infrastructure};

% Column B: lateral branches from PythonEnvironments
\node[fill=colorDesirable,
      right=of pyenv]                       (latex){Ch.6 \LaTeX{}};

% Data challenge track (left column)
\node[fill=colorRequired,
      left=of numop]                        (dc)   {Ch.9 Data Challenge};
\node[fill=colorRequired,
      below=of dc]                          (anal) {Ch.10 Analysis Activities};

% Assessment (bottom, spanning both tracks)
\node[fill=colorRequired,
      below=2.0cm of llm,
      minimum width=9.0cm]                  (asm)  {Ch.11 Assessment \& Progress Record};

% ---- Arrows: main spine (hard dependencies) ----
\draw[harrow] (env)   -- (cli);
\draw[harrow] (cli)   -- (prog);
\draw[harrow] (prog)  -- (pyenv);
\draw[harrow] (pyenv) -- (numop);
\draw[harrow] (numop) -- (git);
\draw[harrow] (git)   -- (llm);

% ---- Arrows: lateral branches ----
\draw[darrow] (pyenv) -- (latex);     % recommended after environments
\draw[harrow] (pyenv.west) to[out=200,in=90] (dc.north);  % env needed for data challenge

% ---- Arrows: data challenge track ----
\draw[harrow] (dc) -- (anal);

% ---- Arrows: Assessment references everything ----
\draw[darrow] (llm)  -- (asm);
\draw[darrow] (anal.south) |- (asm.west);

\end{tikzpicture}

\caption{Chapter dependency structure. Solid arrows indicate hard
dependencies; dashed arrows indicate recommended order. Node color
encodes category: \textcolor{colorRequired}{\textbf{Required}},
\textcolor{colorEssential}{\textbf{Essential}},
\textcolor{colorDesirable}{\textbf{Desirable}},
\textcolor{colorOptional}{\textbf{Optional}}.}
\label{fig:dependencies}
\end{figure}


\begin{landscape}
%================================================================
\section*{Which Chapters Do You Need?}
%================================================================

Use Table~\ref{tab:routing} to identify your reading path based on
your prior experience. Chapters marked \textbf{Required} must be
completed; \textbf{Skim} means read the section headings and run the
exercises but skip the explanatory prose; \textbf{Skip} means the
chapter is optional for you.

\begin{table}[htbp]
\centering
\caption{Recommended reading path by prior experience.}
\label{tab:routing}
\footnotesize
\begin{tabular}{l*{10}{c}}
\toprule
\textbf{Profile}
  & \rotatebox{70}{Env.\ Setup}
  & \rotatebox{70}{CLI}
  & \rotatebox{70}{Intro Prog.}
  & \rotatebox{70}{Python Envs.}
  & \rotatebox{70}{Numerical Ops.}
  & \rotatebox{70}{\LaTeX{}}
  & \rotatebox{70}{Git}
  & \rotatebox{70}{LLM}
  & \rotatebox{70}{Data Challenge}
  & \rotatebox{70}{Analysis} \\
\midrule
New to programming
  & $\bullet$ & $\bullet$ & $\bullet$ & $\bullet$ & $\bullet$
  & $\bullet$ & $\bullet$ & S* & $\bullet$ & $\bullet$ \\
Know Python, new to Linux/CLI
  & $\bullet$ & $\bullet$ & S & $\bullet$ & S
  & $\bullet$ & $\bullet$ & S* & $\bullet$ & $\bullet$ \\
Experienced programmer
  & $\bullet$ & S & -- & $\bullet$ & --
  & $\bullet$ & S & S* & $\bullet$ & $\bullet$ \\
\bottomrule
\end{tabular}
\end{table}

\noindent\footnotesize $\bullet$ = Required \quad S = Skim \quad
-- = Skip \quad
*Chapter~\ref{ch:ollama-infrastructure}: all profiles complete
Section~\ref{sec:apikeys} only; infrastructure sections are optional.
\normalsize


%================================================================
\section*{Progress Record}
%================================================================

The final chapter, Chapter~\ref{ch:assessment} (Assessment and
Progress Record), is a fillable PDF form. It organizes the key
completion items from all chapters into four categories: Required,
Essential, Desirable, and Optional. Open it in Adobe Acrobat Reader
to fill and save your responses. Refer to it throughout the program
as a checklist for readiness.


\newpage
%================================================================
\section*{Estimated Completion Times}
%================================================================

Table~\ref{tab:times} shows estimated completion times per chapter.
Times assume no prior exposure to the topic; students with background
in a given area will complete it faster.

\begin{table}[htbp]
\centering
\caption{Estimated completion time per chapter.}
\label{tab:times}
\begin{tabular}{clr}
\toprule
\textbf{Chapter} & \textbf{Title} & \textbf{Estimated Time} \\
\midrule
1  & Environment Setup & 2 hours \\
2  & Command-Line Interfaces and Shell Scripting & 2 hours \\
3  & Introduction to Programming Environments & 2 hours \\
4  & Python Virtual Environments and Workflow Automation & 1.5 hours \\
5  & Numerical Operations with NumPy & 2 hours \\
6  & Introduction to \LaTeX{} & 1.5 hours \\
7  & Version Control with Git and GitHub & 1.5 hours \\
8  & Local LLM Infrastructure with Ollama & 1 hour (students) / 3 hours (instructors) \\
9  & Data Challenge & 1 hour (reading) \\
10 & Analysis Activities & 4--6 hours \\
11 & Assessment and Progress Record & Ongoing \\
\midrule
   & \textbf{Total (student path)} & \textbf{$\approx$ 19--21 hours} \\
\bottomrule
\end{tabular}
\end{table}

\end{landscape}

